\documentclass[10pt]{mypackage}

\usepackage[uprightgreeks,light]{kpfonts}
\renewcommand{\coloneq}{\coloneqq}
%\renewcommand{\mathbb}{\mathds}
\usepackage{homework}
%\usepackage{notes}

%\usepackage[ backend=bibtex, style = alphabetic, sorting=ynt ]{biblatex}
%\addbibresource{  }

\usepackage{parskip}

\fancyhf{}
\fancyhead[R]{Avinash Iyer}
\fancyhead[L]{Algebraic Topology: Homework 25}
\fancyfoot[C]{\thepage}

\setcounter{secnumdepth}{0}

\begin{document}
\RaggedRight
\section{Revised Problems}%
\begin{problem}[Homework 23, Problem 1]
  Let $\left[ v_0,\dots,v_q \right]$ be a $q$-simplex in $\R^m$ for $m\geq q$. Prove that if $d$ is the diameter of a simplex in the barycentric subdivision of $\left[ v_0,\dots,v_q \right]$, then $d\leq \frac{q}{q+1}\operatorname{diam}\left[ v_0,\dots,v_q \right]$.
\end{problem}
\begin{solution}
  We start by observing that the diameter in the $q$-simplex can be expressed as the maximum distance of the distinct vertices; that is,
  \begin{align*}
    \operatorname{diam}\left( \left[ v_0,\dots,v_q \right] \right) &= \max\set{\left\vert v_i-v_j \right\vert : 0\leq i < j \leq q}.\label{eq:diam_simplex}\tag{$\ast$}
  \end{align*}
  This follows from the fact that if $u$ and $w$ are elements with $\Delta^q$, where $w = \sum_{i=0}^{q}s_iv_i$ has $s_i\geq 0$ and $\sum_{i=0}^{q}s_i = 1$, then
  \begin{align*}
    \left\vert u-w \right\vert &= \left\vert u-\sum_{i=0}^{q}s_iv_i \right\vert\\
                               &= \left\vert \sum_{i=0}^{q}s_i\left( u-v_i \right) \right\vert\\
                               &\leq \sum_{i=0}^{q} s_i \left\vert u-v_i \right\vert\\
                               &\leq \sum_{i=0}^{q} s_i \max\set{\left\vert u-v_i \right\vert : 0\leq i \leq q}\\
                               &= \max\set{\left\vert u-v_i \right\vert : 0\leq i \leq q},
  \end{align*}
  and if $u = \sum_{i=0}^{q}t_iv_i$ with $t_i\geq 0$ and $\sum_{i=0}^{q}t_i = 1$, then for fixed $j$, we have
  \begin{align*}
    \left\vert u-v_j \right\vert &= \left\vert \sum_{i=0}^{q}t_i\left( v_i-v_j \right) \right\vert\\
                                 &\leq \sum_{i=0}^{q} t_i \left\vert v_i-v_j \right\vert\\
                                 &= \sum_{i\neq j} t_i \left\vert v_i-v_j \right\vert\\
                                 &\leq \sum_{i\neq j} t_i \max\set{\left\vert v_i-v_j \right\vert : 0\leq i < j \leq q}\\
                                 &\leq \sum_{i=0}^{q}t_i \max\set{ \left\vert v_i-v_j \right\vert : 0\leq i < j \leq q }\\
                                 &= \max\set{\left\vert v_i-v_j \right\vert : 0\leq i < j \leq q}.
  \end{align*}
  Since this maximum is attained, it is in fact our diameter.

  We start by establishing a bound on the distance of the barycenter to any vertex in our simplex. This gives us
  \begin{align*}
    \left\vert b-v_j \right\vert &= \left\vert \frac{1}{q+1} \sum_{i=0}^{q}\left( v_i-v_j \right) \right\vert\\
                                 &\leq \frac{1}{q+1} \sum_{i=0}^{q} \left\vert v_i-v_j \right\vert\\
                                 &\leq \frac{1}{q+1}\sum_{i\neq j} \max\set{\left\vert v_i-v_j \right\vert : 1\leq i < j \leq q}\\
                                 &= \frac{q}{q+1}\operatorname{diam}\left[ v_0,\dots,v_q \right].
  \end{align*}
  Given some simplex in the barycentric subdivision, its vertices are either vertices in the original simplex, the barycenter of the original simplex, or a barycenter of a face of the simplex, which itself is a $\left( q-1 \right)$-simplex. Labeling a face barycenter as $b'$, we observe that its distance to any of the $v_1,\dots,v_q$ in its face $\sigma_p$ is bounded above by $\frac{q-1}{q}\operatorname{diam}\left( \sigma_p \right)$; the diameter of $\sigma_p$ itself is smaller than the diameter of the entire $\Delta^q$, and $\frac{q-1}{q}\leq \frac{q}{q+1}$ for any $q > 0$. 

  Therefore, the diameter of any simplex in the barycentric subdivision is bounded above by the maximum distance of the barycenter to any vertex in the original simplex, giving the bound of $\frac{q}{q+1}\operatorname{diam}\left[ v_0,\dots,v_q \right]$.
\end{solution}
\section{Current Problems}%
\begin{problem}[Problem 1]
  Let $f\colon \left( X,A \right)\rightarrow \left( Y,B \right)$ be a map such that $f\colon X\rightarrow Y$ and the restriction $A\rightarrow B$ are homotopy equivalences.
  \begin{enumerate}[(a)]
    \item Prove that $f_{\ast}\colon H_q\left( X,A \right)\rightarrow H_q\left( Y,B \right)$ is an isomorphism for all $q\geq 0$.
    \item Prove that the inclusion $f\colon \left( D^n,S^{n-1} \right)\rightarrow \left( D^n,D^n\setminus \set{0} \right)$ is not a homotopy equivalence of pairs, in that there is no map $g$ such that $g\circ f$ and $f\circ g$ are homotopic to identity maps via maps of pairs.
  \end{enumerate}
\end{problem}
\begin{solution}\hfill
  \begin{enumerate}[(a)]
    \item Passing to the long exact sequence gives the following diagram.
      \begin{center}
        % https://tikzcd.yichuanshen.de/#N4Igdg9gJgpgziAXAbVABwnAlgFyxMJZABgBpiBdUkANwEMAbAVxiRAB12BjKCHBAL6l0mXPkIoAjOSq1GLNgAkA+gEcAFAEEAlCCEjseAkQBMM6vWatEIFRoAau-SAyHxRAMzm5VpWvX2pDp6wi6iRhLIACzelgo2KsCqALSSAlpOoa5ixigArLHy1hzcvPwhBjmRZJKyccWcPHyCWeHuUqS1FkV+GgBCmZURpp11PQn+AJqDYW65yF5dPvG2U6QDFbNVRDFL9X5JqekbztnD+aPdvjaNZYKyMFAA5vBEoABmAE4QALZI0iAcBAkAB2K4rd7KYCcOhwHACTZfX5IMyA4GIAAc4OKkOh7Fh8MR3z+iC8aKQAE5sWxcTC4QjnEiSTFyYhJGRljioXTCdQGHQAEYwBgABTauRADBg7xwID5WDAxTgEAYWCgcpAAAsYHR1YgwEwGAxqDg6FgGGxIIqici2Sb0aj9jYsNz8fSbSTUUCkGSnSAAFaugkM0JMn32pAsv2cNB0T54RgeyMRxB5RnEpAc71s9O2gBsKZBuZJYNZWM5bBdeODScxKapFZsger7uLlJT7OpN3YsfjWETbbZWfRkjSFAEQA
\begin{tikzcd}
\cdots \arrow[r] & H_q(A) \arrow[d, "f_{\ast}"] \arrow[r, "i_{\ast}"] & H_q(X) \arrow[d, "f_{\ast}"] \arrow[r, "j_{\ast}"] & {H_q(X,A)} \arrow[d, "f_{\ast}"] \arrow[r, "\partial"] & H_{q-1}(A) \arrow[d, "f_{\ast}"] \arrow[r] & \cdots \\
\cdots \arrow[r] & H_q(B) \arrow[r, "i_{\ast}'"]                       & H_q(Y) \arrow[r, "j_{\ast}'"]                       & {H_q(Y,B)} \arrow[r, "\partial'"]                       & H_{q-1}(B) \arrow[r]                       & \cdots
\end{tikzcd}
      \end{center}
      Here, we have $f_{\ast}$ at $H_q(A)$ and $H_{q}(X)$ are isomorphism, as is the case at $H_{q-1}\left( A \right)$ and $H_{q-1}\left( X \right)$, the latter not pictured.

      Our goal is to show that the middle $f_{\ast}$ is an isomorphism, which is an application of the Five Lemma. However, we have not proven the Five Lemma in class, so we will prove it here.

      First, let $z\in \ker\left( f_{\ast} \right)\subseteq H_q\left( X,A \right)$ be a (representative of a) homology class. Then, we have $f_{\ast}\left( z \right) = 0$, meaning that $\partial'f_{\ast}\left( z \right) = 0$. By commutativity, this means $f_{\ast}\partial \left( z \right) = 0$, meaning that $\partial\left( z \right)\in \ker\left( f_{\ast} \right) = \set{0}$. By exactness, this means we have some $z'\in H_q\left( X \right)$ such that $j_{\ast}\left( z' \right) = z$. By commutativity, we have
      \begin{align*}
        0 &= f_{\ast}\left( z \right)\\
          &= f_{\ast}j_{\ast}\left( z' \right)\\
          &= j_{\ast}'f_{\ast}\left( z' \right),
      \end{align*}
      meaning that $f_{\ast}\left( z' \right)$ is contained in the kernel of $j_{\ast}$. By exactness, there is some $z''\in H_q\left( B \right)$ such that $i_{\ast}'\left( z'' \right) = f_{\ast}\left( z' \right)$. Since $f_{\ast}$ on $H_{q}\left( A \right)$ is surjective, it follows that there is some $z'''\in H_q\left( A \right)$ such that $f_{\ast}\left( z''' \right) = z''$. We have
      \begin{align*}
        f_{\ast}\left( z' \right) &= i'_{\ast}\left( z'' \right)\\
                                  &= i'_{\ast}f_{\ast}\left( z''' \right)\\
                                  &= f_{\ast}i_{\ast}\left( z''' \right).
      \end{align*}
      Since $f_{\ast}$ is an isomorphism on $H_q\left( X \right)$, it follows that $z' = i_{\ast}\left( z'' \right)$. Since $j_{\ast}\left( z' \right) = z$, we have $z = j_{\ast}i_{\ast}\left( z'' \right) = 0$, so $\ker\left( f_{\ast} \right)\subseteq H_q\left( X,A \right)$ is thus equal to zero, and $f_{\ast}$ is thus injective.
  \end{enumerate}
\end{solution}
\end{document}
